\hypertarget{ux437ux430ux434ux430ux43dux438ux435-ux43dux430-ux432ux44bux43fux443ux441ux43aux43dux443ux44e-ux43aux432ux430ux43bux438ux444ux438ux43aux430ux446ux438ux43eux43dux43dux443ux44e-ux440ux430ux431ux43eux442ux443-ux441ux442ux430ux440ux442ux430ux43f}{%
\section{Задание на выпускную квалификационную работу (стартап)}\label{ux437ux430ux434ux430ux43dux438ux435-ux43dux430-ux432ux44bux43fux443ux441ux43aux43dux443ux44e-ux43aux432ux430ux43bux438ux444ux438ux43aux430ux446ux438ux43eux43dux43dux443ux44e-ux440ux430ux431ux43eux442ux443-ux441ux442ux430ux440ux442ux430ux43f}}

\begin{quote}
Актуализация 22.04.2026: формулировки задания синхронизированы с текущим состоянием кода (выполнены промпты 13--22; validation report + PLR/packet-loss обновления).
\end{quote}

Формальный текст задания вынесен из \href{roadmap_task.md}{\texttt{roadmap\_task.md}}, чтобы календарный план и официальное задание не смешивались в одном файле.

\textbf{Согласование с ТЗ на ПО:} детали реализации платформы --- в \href{ТЗ.md}{\texttt{ТЗ.md}}. Раздел \textbf{§10} описывает \textbf{целевую} модель динамического определения угрозы по нескольким источникам (GNSS, IMU, радиоканал, согласованность роя и др.) и слияния признаков. В репозитории реализована \textbf{упрощённая} цепочка по MAVLink (эвристики, без ML/SDR) --- см. \href{Промпты.md}{\texttt{Промпты.md}} 9--12; \textbf{полная} постановка §10 \textbf{не входит} в минимальный объём подтверждения задания и может излагаться в перспективах ВКР и в \href{knowledge/}{\texttt{knowledge/}}.

\begin{center}\rule{0.5\linewidth}{0.5pt}\end{center}

Министерство науки и высшего образования Российской Федерации Федеральное государственное автономное образовательное учреждение высшего образования «СЕВЕРО-КАВКАЗСКИЙ ФЕДЕРАЛЬНЫЙ УНИВЕРСИТЕТ»

\textbf{Факультет:} Факультет математики и компьютерных наук им. профессора Н.И. Червякова\\
\textbf{Кафедра:} Организации и технологии защиты информации\\
\textbf{Направление подготовки:} 10.03.01 «Информационная безопасность»\\
\textbf{Направленность (профиль):} «Организация и технологии защиты информации (по отрасли или в сфере профессиональной деятельности)»

«УТВЕРЖДАЮ» Зав. кафедрой организации и технологии защиты информации канд. техн. наук, доцент Петренко В. И. \_\_\_\_\_\_\_\_\_\_\_\_\_\_\_\_\_\_\_\_\_\_\_\_\_\_\_\_\_ (подпись) «\_\_\_\_\_»\_\_\_\_\_\_\_\_\_\_\_\_\_2026 г.

\textbf{Тема ВКРС:} Разработка платформы автоматизированного планирования безопасных маршрутов для групп сельскохозяйственных дронов с учётом киберугроз\\
\emph{(Утверждена приказом проректора по учебной работе от «\textbf{\emph{»}}\_\_\_\_\_2026 г. №\_\_\_\_\_\_)}

\begin{center}\rule{0.5\linewidth}{0.5pt}\end{center}

\hypertarget{ux440ux430ux437ux434ux435ux43b-1.-ux43fux440ux43eux435ux43aux442ux438ux440ux43eux432ux430ux43dux438ux435-ux430ux440ux445ux438ux442ux435ux43aux442ux443ux440ux44b-ux438-ux440ux430ux437ux440ux430ux431ux43eux442ux43aux430-ux441ux435ux440ux432ux435ux440ux43dux43eux439-ux447ux430ux441ux442ux438-ux43fux43bux430ux442ux444ux43eux440ux43cux44b-ux43fux43bux430ux43dux438ux440ux43eux432ux430ux43dux438ux44f-ux431ux435ux437ux43eux43fux430ux441ux43dux44bux445-ux43cux430ux440ux448ux440ux443ux442ux43eux432-ux431ux43fux43bux430-ux441-ux43cux43eux434ux443ux43bux435ux43c-ux43eux446ux435ux43dux43aux438-ux43aux438ux431ux435ux440ux443ux433ux440ux43eux437}{%
\subsection{РАЗДЕЛ 1. Проектирование архитектуры и разработка серверной части платформы планирования безопасных маршрутов БПЛА с модулем оценки киберугроз}\label{ux440ux430ux437ux434ux435ux43b-1.-ux43fux440ux43eux435ux43aux442ux438ux440ux43eux432ux430ux43dux438ux435-ux430ux440ux445ux438ux442ux435ux43aux442ux443ux440ux44b-ux438-ux440ux430ux437ux440ux430ux431ux43eux442ux43aux430-ux441ux435ux440ux432ux435ux440ux43dux43eux439-ux447ux430ux441ux442ux438-ux43fux43bux430ux442ux444ux43eux440ux43cux44b-ux43fux43bux430ux43dux438ux440ux43eux432ux430ux43dux438ux44f-ux431ux435ux437ux43eux43fux430ux441ux43dux44bux445-ux43cux430ux440ux448ux440ux443ux442ux43eux432-ux431ux43fux43bux430-ux441-ux43cux43eux434ux443ux43bux435ux43c-ux43eux446ux435ux43dux43aux438-ux43aux438ux431ux435ux440ux443ux433ux440ux43eux437}}

\textbf{Студент:} Сутормин Матвей Павлович (группа: ИНБ-б-о-22-1, Курс: 4)

\hypertarget{ux438ux441ux445ux43eux434ux43dux44bux435-ux434ux430ux43dux43dux44bux435-ux434ux43bux44f-ux43fux440ux43eux435ux43aux442ux438ux440ux43eux432ux430ux43dux438ux44f}{%
\subsubsection{2.1. Исходные данные для проектирования:}\label{ux438ux441ux445ux43eux434ux43dux44bux435-ux434ux430ux43dux43dux44bux435-ux434ux43bux44f-ux43fux440ux43eux435ux43aux442ux438ux440ux43eux432ux430ux43dux438ux44f}}

\begin{itemize}
\tightlist
\item
  Математическая модель интегрального критерия доверенной эффективности.
\item
  Алгоритмы динамического перепланирования маршрутов.
\item
  Метрика оценки деструктивного воздействия на процесс координации и маршрутизации.
\item
  Данные о характеристиках коммерческих агро-БПЛА (автономность, полезная нагрузка, радиус связи).
\item
  Open-source библиотеки: OR-Tools, networkx, FastAPI, PostGIS.
\end{itemize}

\hypertarget{ux441ux442ux440ux443ux43aux442ux443ux440ux430-ux438-ux441ux43eux434ux435ux440ux436ux430ux43dux438ux435-ux43fux435ux440ux432ux43eux433ux43e-ux440ux430ux437ux434ux435ux43bux430-ux432ux43aux440ux441}{%
\subsubsection{2.2. Структура и содержание первого раздела ВКРС:}\label{ux441ux442ux440ux443ux43aux442ux443ux440ux430-ux438-ux441ux43eux434ux435ux440ux436ux430ux43dux438ux435-ux43fux435ux440ux432ux43eux433ux43e-ux440ux430ux437ux434ux435ux43bux430-ux432ux43aux440ux441}}

\begin{longtable}[]{@{}
  >{\raggedright\arraybackslash}p{(\columnwidth - 4\tabcolsep) * \real{0.3077}}
  >{\raggedright\arraybackslash}p{(\columnwidth - 4\tabcolsep) * \real{0.3077}}
  >{\centering\arraybackslash}p{(\columnwidth - 4\tabcolsep) * \real{0.3846}}@{}}
\toprule\noalign{}
\begin{minipage}[b]{\linewidth}\raggedright
Структурный элемент ВКРС / этап ГИА
\end{minipage} & \begin{minipage}[b]{\linewidth}\raggedright
Содержание
\end{minipage} & \begin{minipage}[b]{\linewidth}\centering
Компетенции (10.03.01)
\end{minipage} \\
\midrule\noalign{}
\endhead
\bottomrule\noalign{}
\endlastfoot
Титульный лист, задание & Оформление по регламенту ВКРС СКФУ & --- \\
Содержание & Перечень разделов и подразделов работы & --- \\
Введение & Обоснование актуальности. Анализ рынка и существующих решений. Формулировка проблемы, цели и задач. Описание бизнес-идеи стартапа. & ОПК-1, ОПК-3 \\
1.1 Общее описание проекта (совместный раздел) & Тренды рынка агродронов, анализ ЦА, конкурентный анализ. Описание продукта. Бизнес-модель (SaaS). Финансовый план и стратегия. & ОПК-1, ОПК-3, УК-2 \\
1.2 Анализ угроз и моделирование рисков & Систематизация деструктивных воздействий. Разработка модели угроз. Формализация метрики оценки деструктивного воздействия. & ОПК-2, ОПК-4, ПК-1 \\
1.3 Проектирование архитектуры платформы & Разработка общей архитектуры системы. Проектирование базы данных (PostGIS), REST API. Модуль карты рисков. Спецификация MAVLink-интеграции. & ОПК-4, ПК-2, ПК-3 \\
1.4 Разработка серверной части и алгоритмов & Backend (Python, FastAPI). Алгоритм разбиения зоны покрытия. Модуль карты рисков. Логика динамического перепланирования. & ОПК-4, ПК-2, ПК-3, ПК-4 \\
1.5 Эксперименты и оценка эффективности & Тестовые сценарии и метрики по \texttt{План\ обдуманный.md} §1.5. Воспроизводимая Python-симуляция (baseline vs предлагаемый алгоритм); при наличии среды --- демонстрация на ArduPilot SITL. Сравнение с аналогами по ключевым критериям. & ОПК-4, ПК-3, ПК-4 \\
Заключение & Выводы по результатам разработки серверной части и алгоритмов. Перспективы. & УК-1 \\
Список использованных источников & Не менее 40 источников, в т.ч. не менее 10 на иностранном языке & --- \\
Защита ВКРС & Доклад, демонстрация MVP, результаты тестирования & УК-4 \\
\end{longtable}

\hypertarget{ux43fux435ux440ux435ux447ux435ux43dux44c-ux433ux440ux430ux444ux438ux447ux435ux441ux43aux43eux433ux43e-ux43cux430ux442ux435ux440ux438ux430ux43bux430}{%
\subsubsection{2.3. Перечень графического материала:}\label{ux43fux435ux440ux435ux447ux435ux43dux44c-ux433ux440ux430ux444ux438ux447ux435ux441ux43aux43eux433ux43e-ux43cux430ux442ux435ux440ux438ux430ux43bux430}}

Диаграмма архитектуры платформы (UML); ER-диаграмма базы данных; блок-схема алгоритма маршрутизации с учётом риска; блок-схема алгоритма динамического перепланирования; скриншоты работы веб-платформы и/или ArduPilot SITL (по наличию); графики сравнительной оценки эффективности (Python-симуляция).

\textbf{Дата выдачи задания:} 01 апреля 2026 г.\\
\textbf{Руководитель (координатор) проекта:} \_\_\_\_\_\_\_\_\_\_\_\_\_\_ Петренко В.И.\\
\textbf{Руководитель первого раздела проекта:} \_\_\_\_\_\_\_\_\_\_\_\_\_\_ Петренко В.И.\\
\textbf{Задание принял к исполнению:} \_\_\_\_\_\_\_\_\_\_\_\_\_\_ Сутормин М.П.

\begin{center}\rule{0.5\linewidth}{0.5pt}\end{center}

\hypertarget{ux440ux430ux437ux434ux435ux43b-2.-ux440ux430ux437ux440ux430ux431ux43eux442ux43aux430-ux43aux43bux438ux435ux43dux442ux441ux43aux43eux439-ux447ux430ux441ux442ux438-ux43fux43bux430ux442ux444ux43eux440ux43cux44b-ux438-ux43cux43eux434ux443ux43bux44f-ux43eux431ux435ux441ux43fux435ux447ux435ux43dux438ux44f-ux438ux43dux444ux43eux440ux43cux430ux446ux438ux43eux43dux43dux43eux439-ux431ux435ux437ux43eux43fux430ux441ux43dux43eux441ux442ux438-ux432ux437ux430ux438ux43cux43eux434ux435ux439ux441ux442ux432ux438ux44f-ux441-ux433ux440ux443ux43fux43fux43eux439-ux431ux43fux43bux430}{%
\subsection{РАЗДЕЛ 2. Разработка клиентской части платформы и модуля обеспечения информационной безопасности взаимодействия с группой БПЛА}\label{ux440ux430ux437ux434ux435ux43b-2.-ux440ux430ux437ux440ux430ux431ux43eux442ux43aux430-ux43aux43bux438ux435ux43dux442ux441ux43aux43eux439-ux447ux430ux441ux442ux438-ux43fux43bux430ux442ux444ux43eux440ux43cux44b-ux438-ux43cux43eux434ux443ux43bux44f-ux43eux431ux435ux441ux43fux435ux447ux435ux43dux438ux44f-ux438ux43dux444ux43eux440ux43cux430ux446ux438ux43eux43dux43dux43eux439-ux431ux435ux437ux43eux43fux430ux441ux43dux43eux441ux442ux438-ux432ux437ux430ux438ux43cux43eux434ux435ux439ux441ux442ux432ux438ux44f-ux441-ux433ux440ux443ux43fux43fux43eux439-ux431ux43fux43bux430}}

\textbf{Студент:} Яковенко Вячеслав Станиславович (группа: ИНБ-б-о-22-1, Курс: 4)

\hypertarget{ux438ux441ux445ux43eux434ux43dux44bux435-ux434ux430ux43dux43dux44bux435-ux434ux43bux44f-ux43fux440ux43eux435ux43aux442ux438ux440ux43eux432ux430ux43dux438ux44f-1}{%
\subsubsection{3.1. Исходные данные для проектирования:}\label{ux438ux441ux445ux43eux434ux43dux44bux435-ux434ux430ux43dux43dux44bux435-ux434ux43bux44f-ux43fux440ux43eux435ux43aux442ux438ux440ux43eux432ux430ux43dux438ux44f-1}}

\begin{itemize}
\tightlist
\item
  Архитектура серверной части платформы и REST API (результат Раздела 1, Сутормин М.П.).
\item
  Модель угроз из базового НИР (рук. доцент В.И. Петренко). Протоколы взаимодействия с БПЛА (MAVLink).
\item
  Стандарты ИБ: ГОСТ Р 56939-2016, OWASP Top-10, методика ФСТЭК по оценке угроз (2021).
\item
  Библиотеки и фреймворки: React, Leaflet, Vite, Zustand; JWT (реализация на backend); при необходимости обоснования --- TLS/mTLS, cryptography (как направления усиления).
\end{itemize}

\hypertarget{ux441ux442ux440ux443ux43aux442ux443ux440ux430-ux438-ux441ux43eux434ux435ux440ux436ux430ux43dux438ux435-ux432ux442ux43eux440ux43eux433ux43e-ux440ux430ux437ux434ux435ux43bux430-ux432ux43aux440ux441}{%
\subsubsection{3.2. Структура и содержание второго раздела ВКРС:}\label{ux441ux442ux440ux443ux43aux442ux443ux440ux430-ux438-ux441ux43eux434ux435ux440ux436ux430ux43dux438ux435-ux432ux442ux43eux440ux43eux433ux43e-ux440ux430ux437ux434ux435ux43bux430-ux432ux43aux440ux441}}

\begin{longtable}[]{@{}
  >{\raggedright\arraybackslash}p{(\columnwidth - 4\tabcolsep) * \real{0.3077}}
  >{\raggedright\arraybackslash}p{(\columnwidth - 4\tabcolsep) * \real{0.3077}}
  >{\centering\arraybackslash}p{(\columnwidth - 4\tabcolsep) * \real{0.3846}}@{}}
\toprule\noalign{}
\begin{minipage}[b]{\linewidth}\raggedright
Структурный элемент ВКРС / этап ГИА
\end{minipage} & \begin{minipage}[b]{\linewidth}\raggedright
Содержание
\end{minipage} & \begin{minipage}[b]{\linewidth}\centering
Компетенции (10.03.01)
\end{minipage} \\
\midrule\noalign{}
\endhead
\bottomrule\noalign{}
\endlastfoot
Титульный лист, задание & Оформление по регламенту ВКРС СКФУ & --- \\
Содержание & Перечень разделов и подразделов работы & --- \\
Введение & Обоснование актуальности обеспечения ИБ. Анализ угроз каналу управления. Цель и задачи раздела. Роль клиентской части и модуля ИБ. & ОПК-1, ОПК-3 \\
2.1 Общее описание проекта (совместный раздел) & Тренды рынка агродронов, анализ ЦА, конкурентный анализ. Описание продукта. Бизнес-модель (SaaS). Финансовый план (Совместно с Суторминым). & ОПК-1, ОПК-3, УК-2 \\
2.2 Анализ требований к ИБ платформы & Анализ угроз (OWASP Top-10, перехват MAVLink, replay-атаки). Формирование требований к защите. Выбор механизмов защиты. & ОПК-2, ОПК-4, ПК-1 \\
2.3 Разработка клиентской части платформы & Проектирование UX/UI оператора (GeoJSON/KML, настройка миссии, карта React.js + Leaflet.js). Панель диспетчеризации. Интеграция с REST API. & ОПК-4, ПК-2, ПК-3 \\
2.4 Разработка модуля информационной безопасности & Подсистема аутентификации (JWT) и ролевая модель. Анализ и обоснование защиты канала MAVLink (в объёме MVP --- pymavlink; TLS/mTLS/HMAC --- как перспектива усиления). Аудит действий (по объёму ВКР). Защита от типовых веб-уязвимостей. & ОПК-2, ОПК-4, ПК-1, ПК-4 \\
2.5 Тестирование безопасности и интеграционное & Разработка методики тестирования. Пентест веб-интерфейса. Моделирование атак на канал связи. Интеграционное тестирование. & ОПК-4, ПК-3, ПК-4 \\
Заключение & Выводы по результатам разработки. Рекомендации по усилению защищённости. & УК-1 \\
Список использованных источников & Не менее 40 источников, в т.ч. не менее 10 на иностранном языке & --- \\
Защита ВКРС & Доклад, демонстрация интерфейса и модуля ИБ, результаты пентеста & УК-4 \\
\end{longtable}

\hypertarget{ux43fux435ux440ux435ux447ux435ux43dux44c-ux433ux440ux430ux444ux438ux447ux435ux441ux43aux43eux433ux43e-ux43cux430ux442ux435ux440ux438ux430ux43bux430-1}{%
\subsubsection{3.3. Перечень графического материала:}\label{ux43fux435ux440ux435ux447ux435ux43dux44c-ux433ux440ux430ux444ux438ux447ux435ux441ux43aux43eux433ux43e-ux43cux430ux442ux435ux440ux438ux430ux43bux430-1}}

Макеты интерфейса (wireframes); скриншоты реализованного веб-интерфейса; диаграмма потоков данных (DFD) с указанием точек защиты; схема аутентификации и авторизации (JWT); схема канала связи с БПЛА (MAVLink; при необходимости --- варианты усиления TLS/mTLS); результаты проверок / пентеста (отчёт или таблица); скриншоты интеграционного тестирования.

\textbf{Срок представления ВКРС к защите:} 01 июня 2026 г.\\
\textbf{Дата выдачи задания:} 01 апреля 2026 г.\\
\textbf{Руководитель (координатор) проекта:} \_\_\_\_\_\_\_\_\_\_\_\_\_\_ Петренко В.И.\\
\textbf{Руководитель второго раздела проекта:} \_\_\_\_\_\_\_\_\_\_\_\_\_\_ Антонов В.В.\\
\textbf{Задание принял к исполнению:} \_\_\_\_\_\_\_\_\_\_\_\_\_\_ Яковенко В.С.

\begin{center}\rule{0.5\linewidth}{0.5pt}\end{center}

\emph{Текст синхронизирован с \texttt{roadmap\_task.md} (вынесен 19.04.2026). Дополнено ссылкой на \texttt{ТЗ.md} §10 (апрель 2026).}
